\chapter{Business intent as constrained selection}
\label{ch:intent}

\section{Separation of concerns}
The model of Chapters~\ref{ch:symmetry}--\ref{ch:inverse} never sees revenue. It outputs a
ranking of candidate origins and an estimated degradation field. Everything about \emph{which}
remediation runs first when recovery resources are constrained lives in a separate layer with its
own vocabulary. This is the decoupling that a business-intent extension to an intent ontology
provides: the operator declares priority in business terms; the network implementation does not
change.

\section{Revenue at risk}
Let $\mathrm{rev}_u\ge0$ be the revenue density attached to node $u$ (nonzero on cells, zero
elsewhere in our twin; a real deployment attaches it wherever revenue is measured), let
$w_u\ge0$ be an operator-declared profile weight from a small catalogue (mass-event venue, transit
hub, enterprise-critical, residential), and let $\mathrm{reach}(v)$ be the forward closure of $v$
along cause$\to$effect edges, the blast radius.
\begin{equation}
\mathrm{rev},\,w:\V\to\R_{\ge0},\qquad \mathrm{reach}(v)=\{u:\ v\leadsto u\}\cup\{v\}.
\label{eq:revenue}
\end{equation}
The revenue at risk of a candidate origin $v$ is the weighted revenue in its blast radius, scaled by
the estimated degradation there:
\begin{equation}
R(v)=\sum_{u\in\mathrm{reach}(v)}w_u\,\mathrm{rev}_u\,\max\big(0,\hat d_u(t)\big).
\label{eq:rar}
\end{equation}
With $w\equiv1$ this is pure revenue density. With a catalogue of weights and a single-node blast
radius it reduces to the ``profile weight $\times$ revenue density $\times$ service impact''
utility used in industry intent extensions, where $\hat d_u$ plays the role of service impact.
Two candidate origins with identical network signatures but different downstream revenue get
different $R$. That is the entire mechanism behind ``same fault, different response''.

\section{Selection under a budget}
Each candidate $v$ induces a remediation branch with a resource cost $c(v)$ determined by its type
(resetting a cell is cheap, scaling a core function is not). Given the top-$k$ root-cause
candidates and a budget $B$ of recovery resources, the executed order is the greedy solution of
\begin{equation}
\max_{S\subseteq\text{top-}k}\ \sum_{v\in S}R(v)\quad\text{s.t.}\quad\sum_{v\in S}c(v)\le B,
\qquad\text{executed in decreasing order of }R(v)/c(v),
\label{eq:selection}
\end{equation}
with at most one high-risk change per domain per healing cycle (a change-risk concurrency policy);
the remaining branches are held as ranked fallbacks for the verify stage. This is a knapsack with a
ratio-greedy heuristic, deliberately simple so that an operator can read the rule in one line and
change it. The output is serialised as an intent object whose targets
carry domain, action, priority, revenue at risk and resource units, in the shape of a TMF921
intent with a business-priority extension.

\begin{remark}[What the business layer does not do]
It does not move the root-cause ranking. If the model is wrong about the origin, revenue
weighting makes the wrong remediation run first with more conviction. The guard against that is
the counterfactual check of Chapter~\ref{ch:inverse} and the recovery measurement of
Chapter~\ref{ch:closed}, not the intent layer.
\end{remark}
